% called by main.tex

\section*{Abstract}
\addcontentsline{toc}{section}{Abstract}
\label{sec::abstract}

Printymand is an academic web platform for print-on-demand commerce. The project studies how a local marketplace can connect customers, independent designers, printers, and administrators in a single workflow: design discovery, product customization, printer selection, cart management, order placement, tracking, and role-specific dashboards. The front end is implemented with Angular and TypeScript, while the back end is structured with NestJS, PostgreSQL, JWT authentication, role-based authorization, validation, interceptors, and modular controllers. Because the platform is still under development, part of the current user experience relies on mock and hybrid data services while the server-side architecture and database foundations are progressively completed.

This report presents the project context, problem statement, comparative study, requirements, architecture, implementation choices, and validation of the main development cycles. The adopted methodology is the Spiral Model, chosen because the platform combines technical risks, evolving requirements, and multiple user roles. The result is not presented as a production-ready commercial website, but as a coherent engineering PFA prototype that demonstrates the core scope and technical direction of a print-on-demand marketplace.

\newpage
\section*{Keywords}
\addcontentsline{toc}{section}{Keywords}
\label{sec::keywords}

Print-on-demand, marketplace, Angular, NestJS, PostgreSQL, JWT, role-based access control, Spiral Model, software engineering.
